\documentclass{article}
\usepackage[utf8]{inputenc}
\usepackage{graphicx}
\usepackage{svg}
\usepackage{float}
\usepackage{hyperref}
\usepackage{booktabs}
\usepackage{geometry}

\geometry{a4paper, margin=1in}

\title{Study of Deep Q-Network on Atari Games: Pong and Breakout}
\author{Flavien Geoffray}
\date{\today}

\begin{document}

\maketitle

\begin{abstract}
    This report presents a study of the Deep Q-Network (DQN) algorithm applied to the Atari 2600 games Pong and Breakout. We implement and evaluate a DQN agent using the Gymnasium environment. For Pong, a standard architecture was used, while for Breakout, an improved 3-layer convolutional neural network (DQNv2) was implemented to handle the increased complexity. We analyze the training process and performance of the agents through reward and Q-value metrics.
\end{abstract}

\section{Introduction}
Reinforcement Learning (RL) has shown great success in learning control policies directly from high-dimensional sensory inputs. The Deep Q-Network (DQN) algorithm, introduced by Mnih et al., combines Q-learning with deep neural networks to stabilize the learning of value functions. In this study, we apply DQN to two classic Atari games: Pong and Breakout.

\section{Methodology}

\subsection{DQN Algorithm}
The agent interacts with the environment by observing a state $s_t$, taking an action $a_t$, receiving a reward $r_t$, and transitioning to a new state $s_{t+1}$. The goal is to maximize the expected cumulative reward. The Q-function $Q(s, a)$ estimates the expected future reward for taking action $a$ in state $s$. We use a neural network parameterized by $\theta$ to approximate the Q-function.

\subsection{Network Architectures}
We explored two network architectures:
\begin{itemize}
    \item \textbf{Base Architecture (Pong):} A standard convolutional network suitable for simpler dynamics like Pong.
    \item \textbf{Improved Architecture (Breakout - DQNv2):} A deeper network with 3 convolutional layers to capture more complex spatial features required for Breakout.
\end{itemize}

\subsection{Hyperparameters}
The training configurations for the two games are detailed in Table \ref{tab:hyperparams}.

\begin{table}[H]
    \centering
    \begin{tabular}{lcc}
        \toprule
        \textbf{Parameter} & \textbf{Pong} & \textbf{Breakout} \\
        \midrule
        Episodes & 5000 & 10000 \\
        Steps per Episode & 10000 & 8000000 \\
        Memory Capacity & 100000 & 1000000 \\
        Batch Size & 32 & 32 \\
        Learning Rate & 0.00025 & 0.0001 \\
        Gamma ($\gamma$) & 0.99 & 0.99 \\
        Epsilon Start & 1.0 & 1.0 \\
        Epsilon End & 0.1 & 0.1 \\
        Target Update Freq & 10000 & 10000 \\
        Learning Starts & 10000 & 50000 \\
        Optimizer & Adam & Adam \\
        \bottomrule
    \end{tabular}
    \caption{Hyperparameters for Pong and Breakout}
    \label{tab:hyperparams}
\end{table}

\section{Results}

\subsection{Pong}
The agent was trained for 5000 episodes. Figure \ref{fig:pong_reward} shows the evolution of the average reward during training, and Figure \ref{fig:pong_qvalues} displays the mean Q-values.

\begin{figure}[H]
    \centering
    \includesvg[width=0.8\linewidth]{best_run_pong/Reward_Training.svg}
    \caption{Average Reward during Training (Pong)}
    \label{fig:pong_reward}
\end{figure}

\begin{figure}[H]
    \centering
    \includesvg[width=0.8\linewidth]{best_run_pong/Q-Values_Mean.svg}
    \caption{Mean Q-Values during Training (Pong)}
    \label{fig:pong_qvalues}
\end{figure}

The results indicate that the agent successfully learned to play Pong, achieving a high average reward of 19.5.

\subsection{Breakout}
For Breakout, the improved 3-layer CNN model was trained for 10000 episodes. The training progress is illustrated in Figure \ref{fig:breakout_reward} and Figure \ref{fig:breakout_qvalues}.

\begin{figure}[H]
    \centering
    \includesvg[width=0.8\linewidth]{best_run_breakout/Reward_Training.svg}
    \caption{Average Reward during Training (Breakout)}
    \label{fig:breakout_reward}
\end{figure}

\begin{figure}[H]
    \centering
    \includesvg[width=0.8\linewidth]{best_run_breakout/Q-Values_Mean.svg}
    \caption{Mean Q-Values during Training (Breakout)}
    \label{fig:breakout_qvalues}
\end{figure}

The improved architecture allowed the agent to learn effective strategies for Breakout, as evidenced by the increasing reward and Q-values.

\section{Conclusion}
This study demonstrated the effectiveness of DQN on Atari games. While a simpler architecture sufficed for Pong, Breakout benefited from a deeper convolutional network (3 layers) and a larger replay buffer. The results confirm that tailoring the network architecture and hyperparameters is crucial for optimal performance in different environments.

\end{document}
